% ================================================================================
% PEÇA 123 / PAPEL DE TRABALHO PT-13 - CASO SALVADOR: ANÁLISE SOBRE POPULAÇÃO TOTAL
% ================================================================================
% Data: 02/dezembro/2025
% Objetivo: Documentar análise de duplicações em Salvador/BA sobre a população
%           total do CAF (3.175.345 registros), complementando o PT-08 Apêndice I
%           que analisou apenas o censo municipal (1.124 registros).
% Referência: peca123_pt_salvador_populacao_total.pdf
% ================================================================================

\documentclass[11pt,a4paper]{article}

% ========== PACOTES PADRÃO ==========
\usepackage[utf8]{inputenc}
\usepackage[brazil]{babel}
\usepackage[T1]{fontenc}
\usepackage{lmodern}
\usepackage[margin=2.5cm]{geometry}
\usepackage{graphicx}
\usepackage{float}
\usepackage{longtable}
\usepackage{booktabs}
\usepackage{xcolor}
\usepackage{colortbl}
\usepackage{fancyhdr}
\usepackage{amsmath, amssymb}
\usepackage{listings}
\usepackage{enumitem}
\usepackage[hidelinks]{hyperref}
\usepackage[breakable]{tcolorbox}

% ========== CORES TCU OFICIAIS ==========
\definecolor{tcured}{RGB}{204,0,0}      % Vermelho TCU #CC0000
\definecolor{tcublue}{RGB}{0,51,102}    % Azul TCU #003366
\definecolor{tcugreen}{RGB}{0,153,76}   % Verde TCU #006633

% ========== CONFIGURAÇÃO DE CABEÇALHO ==========
\pagestyle{fancy}
\fancyhead[L]{Peça 123 / PT-13 - Caso Salvador (População Total)}
\fancyhead[R]{TCU}
\setlength{\headheight}{14pt}

% ========== CONFIGURAÇÃO DE LISTINGS ==========
\lstset{
    basicstyle=\ttfamily\small,
    frame=single,
    breaklines=true,
    numbers=left,
    numberstyle=\tiny\color{gray},
    keywordstyle=\color{tcublue}\bfseries,
    commentstyle=\color{tcugreen}\itshape,
    stringstyle=\color{tcured},
    columns=fullflexible
}

% ========== INFORMAÇÕES DO DOCUMENTO ==========
\title{Papel de Trabalho de Auditoria Técnica\\
\textbf{Peça 123 / PT-13: Caso Salvador/BA --- Análise de Duplicações sobre População Total}}
\author{Equipe de Auditoria --- Fiscalização CAF-MAPA}
\date{02 de dezembro de 2025}

\begin{document}

\maketitle

% ================================================================================
\section{Objetivo}
% ================================================================================

Este papel de trabalho documenta a análise de duplicações espaciais em Salvador/BA
sobre a \textbf{população total do CAF} (3.175.345 registros ativos), complementando
o PT-08 Apêndice I que analisou apenas o censo municipal (1.124 registros com
precisão $\geq$5 casas decimais).

\textbf{Motivação}: O PT-08 Apêndice I identificou 27 registros duplicados em Salvador,
mas a análise sobre a população total revela uma concentração muito maior de 6.644
imóveis em uma única coordenada.

% ================================================================================
\section{Metodologia}
% ================================================================================

\subsection{Base de Dados}

\begin{itemize}
    \item \textbf{Fonte}: Tabela \texttt{tmp\_populacao\_analisavel}
    \item \textbf{População total}: 3.175.345 registros (CAF ativos em UFs ativas)
    \item \textbf{Filtro}: Código IBGE do município de Salvador (\texttt{cd\_municipio = '2927408'})
    \item \textbf{Critério de duplicação}: Coordenadas exatas (mesmo \texttt{nr\_latitude} e \texttt{nr\_longitude})
\end{itemize}

\subsection{Query SQL Executada}

\begin{lstlisting}[language=SQL, caption={Análise de duplicações em Salvador sobre população total}]
-- Detalhes da coordenada com máxima duplicação em Salvador
SELECT
    nr_latitude,
    nr_longitude,
    COUNT(*) as total_imoveis,
    MIN(dt_criacao) as primeira_criacao,
    MAX(dt_criacao) as ultima_criacao,
    MAX(dt_criacao)::date - MIN(dt_criacao)::date as dias_periodo,
    ROUND(COUNT(*) * 100.0 /
        (SELECT COUNT(*) FROM tmp_populacao_analisavel), 2)
        as percentual_total
FROM tmp_populacao_analisavel
WHERE nr_latitude = -12.960000 AND nr_longitude = -38.510000
GROUP BY nr_latitude, nr_longitude;
\end{lstlisting}

% ================================================================================
\section{Resultados}
% ================================================================================

\subsection{Coordenada com Máxima Duplicação}

\begin{table}[H]
\centering
\caption{Análise da Coordenada Crítica em Salvador/BA}
\label{tab:salvador}
\begin{tabular}{lc}
\toprule
\textbf{Parâmetro} & \textbf{Valor} \\
\midrule
Latitude & -12,960000 \\
Longitude & -38,510000 \\
\midrule
Total de imóveis & \textbf{6.644} \\
Primeira criação & 14/junho/2024 \\
Última criação & 02/setembro/2025 \\
Período total & \textbf{445 dias} \\
\midrule
\% sobre população total & 0,21\% \\
\% sobre duplicações nacionais & 1,46\% \\
\bottomrule
\end{tabular}
\end{table}

\subsection{Comparação com PT-08 Apêndice I}

\begin{table}[H]
\centering
\caption{Comparação entre Análises}
\label{tab:comparacao}
\begin{tabular}{lcc}
\toprule
\textbf{Parâmetro} & \textbf{PT-08 Apêndice I} & \textbf{Peça 123 (Este documento)} \\
\midrule
Base de análise & Censo municipal & População total \\
Registros analisados & 1.124 & 3.175.345 \\
\midrule
Imóveis na coordenada & 27 & 6.644 \\
Período analisado & 21 dias & 445 dias \\
Percentual & 2,40\% (sobre censo) & 1,46\% (sobre duplicações) \\
\bottomrule
\end{tabular}
\end{table}

\textbf{Explicação da diferença}: O PT-08 Apêndice I analisou apenas os 1.124 registros
de Salvador que possuíam precisão $\geq$5 casas decimais, encontrando 27 duplicados.
O presente PT-13 analisa todos os registros de Salvador na população total, incluindo
aqueles com menor precisão, revelando que a coordenada genérica (-12,96; -38,51)
concentra 6.644 cadastros.

% ================================================================================
\section{Contexto Geográfico}
% ================================================================================

A coordenada (-12,960000; -38,510000) corresponde a uma área no \textbf{centro urbano
de Salvador/BA}. A utilização desta coordenada para 6.644 cadastros distintos,
ao longo de 445 dias, evidencia:

\begin{enumerate}
    \item \textbf{Falha de validação semântica}: O sistema aceita coordenadas
    sintaticamente válidas (dentro dos intervalos permitidos), mas semanticamente
    absurdas (área urbana consolidada declarada como zona rural).

    \item \textbf{Persistência do problema}: O intervalo de 445 dias (junho/2024 a
    setembro/2025) demonstra que a vulnerabilidade não foi corrigida ao longo do tempo.

    \item \textbf{Escala do problema}: 6.644 cadastros em uma única coordenada
    representam 1,46\% de todas as duplicações espaciais identificadas no CAF nacional.
\end{enumerate}

% ================================================================================
\section{Memória de Cálculo}
% ================================================================================

\subsection{Cálculo do Percentual sobre Duplicações Nacionais}

\begin{align*}
\text{Total de duplicações nacionais} &= 453.657 \text{ imóveis} \\
\text{Imóveis em Salvador (coord. crítica)} &= 6.644 \text{ imóveis} \\
\text{Percentual} &= \frac{6.644}{453.657} \times 100 = 1,46\%
\end{align*}

\subsection{Query de Verificação}

\begin{lstlisting}[language=SQL, caption={Cálculo do percentual sobre duplicações nacionais}]
WITH todas_duplicacoes AS (
    SELECT nr_latitude, nr_longitude, COUNT(*) as qtd
    FROM tmp_populacao_analisavel
    GROUP BY nr_latitude, nr_longitude
    HAVING COUNT(*) > 1
)
SELECT
    SUM(qtd) as total_duplicados_nacional,
    6644 as imoveis_salvador,
    ROUND(6644 * 100.0 / SUM(qtd), 2) as percentual
FROM todas_duplicacoes;
-- Resultado: 453.657 | 6.644 | 1,46%
\end{lstlisting}

% ================================================================================
\section{Conclusão}
% ================================================================================

A análise sobre a população total do CAF revela que a coordenada (-12,960000; -38,510000)
em Salvador/BA concentra \textbf{6.644 imóveis} cadastrados ao longo de \textbf{445 dias},
representando \textbf{1,46\%} de todas as duplicações espaciais identificadas no sistema.

Este achado complementa o PT-08 Apêndice I e demonstra que:

\begin{itemize}
    \item A magnitude do problema é significativamente maior quando analisada sobre
    a população total (6.644 vs. 27 registros);
    \item A coordenada genérica é utilizada como ``ponto default'' para cadastros
    sem georreferenciamento preciso;
    \item A falha persiste por período prolongado (445 dias), indicando ausência
    de controles preventivos efetivos.
\end{itemize}

% ================================================================================
\section{Anexos}
% ================================================================================

\subsection{Parâmetros de Conexão}

\begin{verbatim}
Host: 172.23.80.1
Port: 5433
Database: caf_mapa
User: postgres
Data de execução: 02/dezembro/2025
\end{verbatim}

\subsection{Artefatos Gerados}

\begin{itemize}
    \item \texttt{peca123\_pt\_salvador\_populacao\_total.tex} --- Este documento (código fonte)
    \item \texttt{peca123\_pt\_salvador\_populacao\_total.pdf} --- Versão compilada (Peça 123)
\end{itemize}

\end{document}
